\documentclass[11pt,a4paper]{article}
\usepackage[T1]{fontenc}
\usepackage{lmodern,amsmath,amssymb,amsthm,microtype}
\usepackage[margin=27mm]{geometry}
\usepackage[colorlinks=true,linkcolor=blue,urlcolor=blue]{hyperref}
\setlength{\emergencystretch}{2em}
\newtheorem{proposition}{Proposition}
\newcommand{\Q}{\mathbb Q}
\title{A small quadratic admissible trace class}
\author{}\date{17 November 2026}
\begin{document}\maketitle

Retain a native detector weight $w$ and its actual rational trace
$V=\operatorname{Tr}(W^{-1})$. Let
\[
 C=V^2+w(2-V),\qquad \delta(V)=V^2-V+2.
\]
The exceptional set contains two. Outside this set every odd valuation
of the physical weight is positive, as in the normalized general and
conditional sources. All existing point, denominator and detector
guards are retained.

\begin{proposition}
The rational-square condition $C\delta(V)\in\Q^{\times2}$ excludes
every negative odd valuation of $b=4w/V^2$ outside the exceptional set.
At every prepared real point-algebra character the radicand is strictly
positive. The additional class introduces no free coordinate,
denominator, guard, or loss of contact primes in the soundness proof.
\end{proposition}
\begin{proof}
At a hypothetical negative odd valuation of $b$, put
$\nu=v_p(w)$ and $m=v_p(V)$. The weight property gives an odd
$\nu>0$, and $2m>\nu$, so $m\ge1$. Thus $2-V$ is a unit and
$v_p(C)=\nu$. Since $\delta(V)\equiv2\pmod p$, this class is a
unit. The product has odd valuation and cannot be a rational square.
Local squareness of the class is not required at these bad primes.

The identity
\[
 \delta(V)=(V-\tfrac12)^2+\tfrac74
\]
gives strict positivity for every real $V$. Prepared characters have
$w>0$ and $0<V\le2$, whence $C>0$. The class is a polynomial in
the already guarded trace. It has no additional denominator and cannot
vanish at any prepared real character, so its finite-algebra norm is
a unit there. The pole argument loses no prime, and the existing
unconditional Lucas or conditional height proof is unchanged.
\end{proof}

\begin{proposition}
The class is compatible with all selected comparison contacts in the
retained completeness construction and with arbitrary fixed-place
preparation near $V=2$.
\end{proposition}
\begin{proof}
Every selected comparison prime is split in $\Q(\sqrt2)$. At a
successful contact the actual trace has residue one or two, while
$w$ has positive valuation. Hence $C$ has residue one or four and
$\delta(V)$ has residue two or four. Both residues are nonzero
squares because two is a square at a split odd prime. Hensel's lemma
applies. This uses the split-contact condition from the retained prime
pattern, not a newly imposed splitting requirement.

At every fixed finite place, $V$ can be prepared arbitrarily close to
two and $\delta(2)=4$. The quotient $\delta(V)/4$ therefore lies
in the local square neighborhood for sufficient preparation. At two
one may explicitly write $V=2+e$ and use
\[
 \delta(2+e)=4+3e+e^2;
\]
the sufficient condition $v_2(e)\ge5$ gives
$\delta(V)/4\in1+8\mathbb Z_2$. The preparation of $C$ itself
is unchanged. The same argument treats primes dividing the fixed
scalar or earlier denominators by increasing their already prescribed
finite preparation depths.
\end{proof}

The exact polynomial identity is
\[
 C\delta(V)=V^4-V^3-wV^3+2V^2+3wV^2-4wV+4w.
\]
It has total degree four, seven monomials and coefficient height four
as a polynomial in $(V,w)$. These are radicand statistics, not degree
or height bounds for a complete defining polynomial. Finite disjunction
with the other admissible classes retains the same external
coordinates and the same one-coordinate encoder. A twenty-one- or
nine-coordinate definition still requires simultaneous rational
realization on the actual detector traces; no such realization is
asserted here.

\section*{Exact retained source references}
The normalized general completeness proof, in the subsection on
comparison primes immediately after its extra-divisibility condition,
states that the selected prime is split in $K$; see
\cite[lines 811--835]{General}. The conditional completeness proof
uses the same inherited prime pattern and its two conjugate
annihilating-order families; see \cite[lines 559--595]{Conditional}.
The explicit split-place statement and the fact that two is a local
square are also isolated in \cite[Lemma ``Fresh split contacts'',
lines 382--404]{Native}. These line numbers refer to the versioned
source filenames below, not to a moving unversioned draft.

\begin{thebibliography}{9}\raggedright
\bibitem{General}\emph{A defining polynomial in twenty-five variables:
combined scalar and comparison conditions}, repository manuscript,
\nolinkurl{papers/few_variables25/few_variables25_merged_20261110_v19.tex}.
\bibitem{Conditional}\emph{Thirteen Variables from Direct Elliptic
Index Comparison}, repository manuscript,
\nolinkurl{papers/conditional_few_variables/direct_input13_20261110_v22.tex}.
\bibitem{Native}\emph{A guarded rational-square criterion for normalized
elliptic detectors}, repository manuscript,
\nolinkurl{papers/pole_parameters/native_square_criterion_20261116_v53.tex}.
\end{thebibliography}
\end{document}
